% Related work (draft). Sources: docs/novelty_check.md (W1). \todo{} marks citations still to be found;
% refs marked VERIFY in refs.bib were read from abstracts only.
\section{Related work}
\label{sec:related}

\subsection{Offloading dependent tasks}
Offloading in mobile edge computing was first studied for independent tasks, where each task is placed on its own and
the objectives add up over tasks. Applications with dependent tasks need a scheduling view: the makespan follows the
critical path of the task graph, and the data passed between tasks cost time and energy whenever the two tasks run on
different nodes. List-scheduling heuristics such as HEFT \citep{topcuoglu2002heft} build a schedule task by task in a
priority order, and many multi-objective evolutionary algorithms optimise placement and order for crisp DAG
offloading in edge or cloud--edge--end systems \todo{cite 2--3 crisp DAG-offloading MOEAs}.
% Candidates in docs/novelty_check.md: the FGCS 2023 MOEA/D cloud-edge-end paper, the FGCS 2025 IIoT dependency-aware paper.
These models treat bandwidth, CPU speed and workload as known values.
% VERIFY against the full text of Shi et al. (only the abstract was read): migration rather than placement, no cloud tier.
\citet{shi2023workflowmigration} consider
dependent workflows with interval bandwidth and computing power, but for workflow migration between edge servers
rather than placement and ordering over devices, edge servers and the cloud. Our model is the first, to our knowledge,
to compute interval makespan along the critical path of a DAG over all three tiers, with both placement and order as
decision variables, and it reduces to the independent-task model of \citet{jin2026imoeadqn} when every
application has one task.

\subsection{Interval multi-objective evolutionary algorithms}
Interval multi-objective problems are solved either by turning them into crisp problems, for example through interval
midpoints and widths, or by comparing intervals directly. IP-MOEA \citep{limbourg2005ipmoea} extends Pareto dominance to
interval vectors and crowds by hypervolume contributions of the interval bounds. Decomposition-based methods compare
solutions on scalarised subproblems: IMOEA/D-AWN \citep{jin2023imoeadawn} adapts the weight vectors and neighbourhood sizes
using an interval sparsity measure, and IMOEA/D-DAW \citep{zhang2026imoeaddaw} uses interval Tchebycheff and penalty-based
scalarisations with dynamic neighbourhoods for robot path planning. IMEA-IED \citep{cui2024imeaied} compares intervals by
an information-entropy dominance and selects within niches. IMaOEACI \citep{cai2025imaoeaci} steers an interval
many-objective search by an individual comprehensive indicator for vehicle selection in federated learning. IMOEA-DQN \citep{jin2026imoeadqn} is the method closest to ours: a three-tier
environmental selection by interval dominance degree, interval crowding distance and uncertainty dominance, with the
shares chosen by a deep Q-network \citep{mnih2015dqn}. None of these methods uses a weight-vector rank as one tier of a
dominance-, diversity- and uncertainty-aware selection.

\subsection{Combining dominance and decomposition}
The loss of selection pressure when most of a population is non-dominated is well known for problems with many
objectives \citep{ishibuchi2008review}; we show that it already occurs with three interval objectives under
possibility-degree dominance. Several crisp algorithms combine Pareto dominance with weight vectors or reference
points: NSGA-III \citep{deb2014nsga3} niches the last front by reference directions, MOEA/DD \citep{li2015moeadd} uses
dominance and decomposition together, and $\theta$-DEA \citep{yuan2016thetadea} defines a dominance relation on
decomposition distances. The rank \ours{} uses comes from MSOPS \citep{hughes2003msops}, which ranks a solution by its best
position over the orderings induced by a set of weighted min--max scores; Global WASF-GA \citep{saborido2017gwasfga}
classifies a population into fronts by achievement scalarising functions over weight vectors. \ours{} applies the
MSOPS rank to interval midpoints and uses it as the convergence tier inside the three-tier selection, which keeps the
diversity and uncertainty tiers of the interval setting.

\subsection{Learned control of evolutionary algorithms}
Reinforcement learning has been used to choose operators, parameters or selection weights of evolutionary algorithms
online \todo{cite a survey of RL-assisted EAs}. IMOEA-DQN \citep{jin2026imoeadqn} learns the selection
weights from the change in a population indicator after every generation. Negative results are rarely reported; our
analysis in Section~\ref{sec:learned} shows that for this problem the per-generation reward is too noisy, and the room
for an adaptive schedule too small, for learned control to beat a fixed action, and that a graph-attention
encoder \citep{velickovic2018gat} of the DAG does not change this.
